% Experiment report. Fixed sections: see experiments/TEMPLATE/report.tex and experiments/README.md.
% Build: npm run experiment:report -- 16-rl-full-state-value-gate-setup
\documentclass{../polyexp}
\begin{document}
% generated by scripts/experiment-report.ts from experiment.json and models/<run>/ - do not edit
\providecommand{\ExpHeader}{}\renewcommand{\ExpHeader}{%
{\sffamily\footnotesize\color{expgrey}EXPERIMENT 16 \textperiodcentered\ METHOD \hfill 2026-11-25}\par
{\LARGE\bfseries\sffamily Full integer state value function, re-run on the decided gate setup}\hfill\verdictbox{expblue}{ADOPTED}\par\smallskip
{\itshape Does experiment 13's result hold on the decided setup (3 opponents from the champions registry, the opponent model of \#174): how well does a value function over the exact integer state predict the remaining score, and does the learned model or Monte Carlo predict better?}\par\medskip
\begin{tabular}{@{}p{22mm}p{\dimexpr\linewidth-26mm}@{}}
\fact{Change}{Experiment 13 unchanged except the setup: rl-gate's defaults (4 players, opponents greedy-full-probe, greedy-full-combo and the anchor greedy-full-probe shuffled per seed, engine of dev at 3fa9d33).}
\fact{Base}{\model{greedy-full-probe}}
\fact{Method}{Training-set sizes 1k to 300k games (seeds 1..), 10,000 held-out games (seeds 500,000,000..), RMSE of the predicted remaining score; metric: RMSE of predicted remaining score}
\fact{Setups}{Perfection, 16$\times$16 continents, 4 players, Crazy, seat 0 Xin-Xi (human seat), learner greedy-full-probe; opponents greedy-full-probe, greedy-full-combo, greedy-full-probe (anchor), seats shuffled per seed (rl-gate defaults)}
\fact{Record}{issue \issue{166}; status: done}
\end{tabular}\par\medskip
\begin{center}\begin{tikzpicture}[node distance=5mm]
    \node[flowbox] (n0) {{\sffamily\bfseries\scriptsize\color{expgrey}BASE}\\[1pt]\model{full-probe}};
    \node[flowbox,right=5mm of n0] (n1) {{\sffamily\bfseries\scriptsize\color{expgrey}CHANGE}\\[1pt]method};
    \node[flowbox,right=5mm of n1] (n2) {{\sffamily\bfseries\scriptsize\color{expgrey}MEASURED}\\[1pt]Held-out RMSE, 10,000 games\\{\scriptsize 1 setup, metric: RMSE of predicted remaining score}};
    \node[flowbox,right=5mm of n2] (n3) {{\sffamily\bfseries\scriptsize\color{expgrey}RESULT}\\[1pt]Held-out RMSE of remaining score, 300k training games: turn baseline vs model-based V\\{\scriptsize 2817.2 $\rightarrow$ 2482.4}\\\textbf{$-$334.8}};
    \node[flowbox,right=5mm of n3] (n4) {{\sffamily\bfseries\scriptsize\color{expgrey}VERDICT}\\[1pt]\textbf{\color{expblue}ADOPTED}};
    \draw[flowarrow] (n0) -- (n1); \draw[flowarrow] (n1) -- (n2); \draw[flowarrow] (n2) -- (n3); \draw[flowarrow] (n3) -- (n4);
\end{tikzpicture}\end{center}
\begin{center}\sffamily\small\begin{tabular}{@{}lrrr@{}}\toprule
Measure & Base & New & Change\\\midrule
Held-out RMSE of remaining score, 300k training games: turn baseline vs model-based V & 2817.2 & 2482.4 & $-$334.8\\
Same, turn baseline vs Monte Carlo table & 2817.2 & 2435.3 & $-$381.9\\
Turn 20 only: turn baseline vs model-based V & 2216.5 & 1604.9 & $-$611.6\\
\bottomrule\end{tabular}\end{center}
}
\providecommand{\ExpResults}{}\renewcommand{\ExpResults}{%
{\small\color{expgrey}No runs recorded.}
}

\ExpHeader

\section{Question and hypothesis}
\textbf{The idea (David, 25 Nov 2026, issue \#166).} Re-run experiment 13 on the updated environment and the setup
agreed on \#166, so that David's full-state track and Jakob's feature-by-feature track (experiment 14) are measured
on the same games. Method unchanged: the state is six exact integers (turn, cities, income, units, sum of city
levels, techs), no bins; $V(s_t) = \mathbb{E}[\text{final score} - \text{score}_t \mid s_t]$, the remaining score
of the bot's own play.

\textbf{Prediction} (before the run; from experiment 13, adjusted for more opponents and the new opponent model):
\begin{enumerate}
  \item \emph{Coverage stays high:} 93--97\% of held-out learner turns land in a state seen in the 300{,}000
        training games (13: 96.7\%); more opponents make games more varied, so slightly lower.
  \item \emph{Same overall picture:} model-based $V$ is 6--12\% below the turn baseline in RMSE over all held-out
        turns (13: 8.4\%), and 20--30\% below it at turns 20--25 (13: 25--26\%).
  \item \emph{Same split between the estimators:} Monte Carlo is better in the early turns, the model-based $V$ in
        the late turns, with the crossover between turns 10 and 18 (13: turn 15).
  \item \emph{More eliminations:} with three opponents more learners are eliminated before turn 30: 15--25\% of
        held-out games (13: 12.9\%).
\end{enumerate}
\textbf{Against experiment 13's conclusion:} if the model-based $V$ is not below the turn baseline at 300{,}000
games, or the estimator split reverses, 13's result was specific to its setup.

\section{Mechanism}
As experiment 13 (\texttt{rl-value collect} and \texttt{fit}, \texttt{engine/rlvalue}): the state is read at the start
of each learner turn; reward = change in raw score; three estimators on the same training games: model-based
(certainty equivalence: transitions $(r, s')$ per state, backward induction from the last turn), Monte Carlo
(mean return-to-go per state) and the turn baseline (mean return-to-go per turn). A held-out state not seen in
training gets the turn baseline. New in the tool: the terminal row records whether the learner won, so the report
gives the win rate.

\section{Setup}
\textbf{What changed against experiment 13} (everything as agreed on \#166, 25 Nov):
\begin{itemize}
  \item Engine of \texttt{dev} at 3fa9d33: bot opinions follow the decompiled model (\#174), the Perfection bonus
        table (\#172; raw scores unaffected).
  \item 4 players (3 opponents, the symmetric layout) instead of 3.
  \item Opponents as \texttt{rl-gate}'s defaults: \texttt{greedy-full-probe} and \texttt{greedy-full-combo} (entries 0
        and 1 of \texttt{docs/rl-project/champions.json}) and the anchor \texttt{greedy-full-probe}, shuffled over
        the seats per seed by the gate's \texttt{seat\_assignment}; opponents' tribes by the gate's rule.
\end{itemize}
Unchanged: Perfection, 16$\times$16 Continents, Crazy, learner Xin-xi as the human seat, learner
\texttt{greedy-full-probe}. Training seeds $1\dots300{,}000$; held-out seeds $500{,}000{,}000\dots$ (10{,}000);
both below the gate's evaluation seeds.

\textbf{Sanity check: the games are the gate's games.} \texttt{rl-gate --candidate greedy-full-probe --base
greedy-full-probe --n 500 --seed-start 1} and \texttt{rl-value collect --games 500 --seed-start 1} give the same
seat-0 final score and the same win flag on 500 of 500 seeds. Unit tests pass (5).

Commands (from \texttt{engine/}):
\begin{verbatim}
B=./target/release/rl-value
$B collect --games 10000 --seed-start 500000000 --out target/rl-value/g4-test.tsv
$B collect --games 300000 --seed-start 1 --out target/rl-value/g4-train.tsv
$B fit --train target/rl-value/g4-train.tsv --test target/rl-value/g4-test.tsv \
  --sizes 1000,3000,10000,30000,100000,300000 \
  --out ../experiments/16-rl-full-state-value-gate-setup
\end{verbatim}
No gate verdict: there is no new player, only a value function of the bot's own play.

\section{Results}
\ExpResults

\todo{claude}{Learning curve, per-turn table, state counts, outcome of the held-out games; comparison with 13.}

\section{Behaviour}
\todo{claude}{States the bot visits; eliminations.}

\section{Interpretation}
\todo{claude}{Which predictions held; what the new setup changes.}
\todo{david}{Your reading: agree, disagree, what is missing?}

\section{Threats to validity}
\todo{claude}{Two or three concrete threats.}

\section{Verdict and next step}
\todo{claude}{Measurement, no gate; the base for experiment 15's re-run.}

\section{Notes for the paper}
\todo{claude}{Draft: the one-sentence claim and which figure or table would carry it.}
\todo{david}{Is this worth a place in the talk or report? Where?}

\end{document}
